\documentclass[11pt]{article}
\usepackage{preamble}
\setlength{\parskip}{4pt}

\begin{document}

\daybanner{AP Statistics \textperiodcentered\ Unit 1 \textperiodcentered\ Topic 1.7 \textperiodcentered\ B: 2026-09-21 \textperiodcentered\ E: 2026-09-25 \textperiodcentered\ TEACHER KEY}{Which Summary Tells the Story?}

\sectionbanner{CED TETHER}

{\small
\begin{itemize}[leftmargin=1.5em,itemsep=2pt,topsep=2pt]
  \item ENDURING UNDERSTANDING: UNC-1 Graphical representations and statistics allow us to identify and represent key features of data.
  \item SKILLS:
  \item \textbf{Skill 2.C} --- Calculate summary statistics, relative positions of points within a distribution, correlation, and predicted response.
  \item \textbf{Skill 4.B} --- Interpret statistical calculations and findings to assign meaning or assess a claim.
  \item LO UNC-1.I Calculate measures of center and position for quantitative data. [Skill 2.C]
  \item EK UNC-1.I.1 A statistic is a numerical summary of sample data.
  \item EK UNC-1.I.2 The mean is the sum of all the data values divided by the number of values. For a sample, the mean is denoted by x-bar: $\bar{x}$ = (1/n) $\sum$ $x_i$, where $x_i$ represents the $i^{\mathrm{th}}$ data point in the sample and n represents the number of data values in the sample.
  \item EK UNC-1.I.3 The median of a data set is the middle value when data are ordered. When the number of data points is even, the median can take on any value between the two middle values. In AP Statistics, the most commonly used value for the median of a data set with an even number of values is the average of the two middle values.
  \item EK UNC-1.I.4 The first quartile, Q1, is the median of the half of the ordered data set from the minimum to the position of the median. The third quartile, Q3, is the median of the half of the ordered data set from the position of the median to the maximum. Q1 and Q3 form the boundaries for the middle 50\% of values in an ordered data set.
  \item EK UNC-1.I.5 The p\textsuperscript{th} percentile is interpreted as the value that has p\% of the data less than or equal to it.
  \item LO UNC-1.J Calculate measures of variability for quantitative data. [Skill 2.C]
  \item EK UNC-1.J.1 Three commonly used measures of variability (or spread) in a distribution are the range, interquartile range, and standard deviation.
  \item EK UNC-1.J.2 The range is defined as the difference between the maximum data value and the minimum data value. The interquartile range (IQR) is defined as the difference between the third and first quartiles: Q3 - Q1. Both the range and the interquartile range are possible ways of measuring variability of the distribution of a quantitative variable.
  \item EK UNC-1.J.3 Standard deviation is a way to measure variability of the distribution of a quantitative variable. For a sample, the standard deviation is denoted by $s_x$: $s_x$ = $\surd$[(1/(n-1)) $\sum$($x_i$ - $\bar{x}$)\textsuperscript{2}]. The square of the sample standard deviation, $s^2$, is called the sample variance.
  \item EK UNC-1.J.4 Changing units of measurement affects the values of the calculated statistics.
  \item LO UNC-1.K Explain the selection of a particular measure of center and/or variability for describing a set of quantitative data. [Skill 4.B]
  \item EK UNC-1.K.1 There are many methods for determining outliers. Two methods frequently used in this course are: i. An outlier is a value greater than 1.5 $\times$ IQR above the third quartile or more than 1.5 $\times$ IQR below the first quartile.
  \item ii. An outlier is a value located 2 or more standard deviations above, or below, the mean.
  \item EK UNC-1.K.2 The mean, standard deviation, and range are considered nonresistant (or non-robust) because they are influenced by outliers. The median and IQR are considered resistant (or robust), because outliers do not greatly (if at all) affect their value.
\end{itemize}
}
\textbf{Essential question:} How can a statistic be calculated correctly but still be the wrong summary to report?\par
\textbf{DOK-3 skill:} 4.B \textperiodcentered\ \textbf{FRQ pattern:} {\small\texttt{whose-\allowbreak{}computation-\allowbreak{}and-\allowbreak{}what-\allowbreak{}it-\allowbreak{}misleads-\allowbreak{}then-\allowbreak{}choose-\allowbreak{}the-\allowbreak{}summary}}\par
\textbf{DOK rationale:} Part (c) diagnoses two methods, selects a summary for purpose, and explains the rejected pair's consequence from the dotplot.\par

\frameworkphaseheader{Do Now $\rightarrow$ Explore $\rightarrow$ Exit}{1 $\rightarrow$ 3}{45}{%
\begin{itemize}[leftmargin=1.5em,itemsep=2pt,topsep=2pt]
  \item Show both recommendations without endorsement; start the video at minute 5.
  \item At minute 33, press calculation versus choice; collect and score part (c), E/P/I.
\end{itemize}
}{%
\begin{itemize}[leftmargin=1.5em,itemsep=2pt,topsep=2pt]
  \item Choose with evidence, follow along, then finish (a)--(c) and the exit.
\end{itemize}
}{%
\begin{itemize}[leftmargin=1.5em,itemsep=2pt,topsep=2pt]
  \item Where does the overall median go when you form the halves?
  \item Can correct arithmetic support a less representative summary?
  \item What changes if the 30-hour point moves farther right?
  \item What could 7.5 suggest without the dotplot?
\end{itemize}
}{Circulate five pairs; press separately on quartiles and summary choice.}

\begin{callout}{\IconWarn\ WATCH FOR}{calloutred}
\begin{itemize}[leftmargin=1.5em,itemsep=2pt,topsep=2pt]
  \item Including the overall median in both halves.
  \item Treating either student as entirely correct.
  \item Choosing resistant summaries but retaining IQR 6.5.
\end{itemize}
\end{callout}

\sectionbanner{SCORING PART (c) --- E / P / I}

\textbf{Expected elements}
\begin{itemize}[leftmargin=1.5em,itemsep=2pt,topsep=2pt]
  \item \textbf{adjudicates-quartiles} (required) --- Identifies Kai's IQR 8 and explains Rosa included the median in both halves
  \item \textbf{selects-resistant-pair} (required) --- Recommends median 8 and corrected IQR 8
  \item \textbf{cites-isolated-thirty} (required) --- Cites the isolated 30, 15-hour gap, or right skew
  \item \textbf{states-nonresistant-consequence} (required) --- Explains how mean and SD could make outlier-driven spread seem typical
\end{itemize}

{\small\renewcommand{\arraystretch}{1.25}
\noindent\begin{tabular}{|p{0.5in}|p{5.9in}|}
\hline
\textbf{E} & Corrects Rosa's method, chooses median and IQR, and links isolated 30 to a mean-and-SD misread. \\ \hline
\textbf{P} & Gets the method or summary choice right but omits a required link. \\ \hline
\textbf{I} & Treats one student as wholly correct or gives no feature-based judgment. \\ \hline
\end{tabular}}

\textbf{Common mistakes}
\begin{itemize}[leftmargin=1.5em,itemsep=2pt,topsep=2pt]
  \item Accepting one report wholesale instead of separating method from summary choice
\end{itemize}

\clearpage
\focusbanner{TODAY'S PROBLEM --- annotated}

Suppose a principal asks, ``How much paid work do students do in a week?'' Fifteen students report these weekly hours: 0, 0, 0, 4, 5, 6, 6, 8, 8, 10, 10, 12, 12, 15, 30.\par\smallskip \textbf{Kai:} ``$\bar{x}=8.4$, $s_x\approx7.5$, and $\mathrm{IQR}=8$. Report the mean and SD because they use every observation.''\par \textbf{Rosa:} ``Median $=8$, $Q_1=4.5$, $Q_3=11$, and $\mathrm{IQR}=6.5$. Report the median and IQR because 30 looks unusual.''\par
\begin{center}
\begin{tikzpicture}
\begin{axis}[
  width=0.68\linewidth,
  height=1.28in,
  ymin=0, ymax=4, ytick=\empty, axis y line=none, axis x line=bottom,
  xlabel={Weekly paid-work hours}, tick label style={font=\small}, label style={font=\small}
]
\addplot[only marks, mark=*, mark size=1.9pt, color=dayblue] coordinates {
  (0,1)
  (0,2)
  (0,3)
  (4,1)
  (5,1)
  (6,1)
  (6,2)
  (8,1)
  (8,2)
  (10,1)
  (10,2)
  (12,1)
  (12,2)
  (15,1)
  (30,1)
};
\end{axis}
\end{tikzpicture}
\end{center}
\begin{firsttakebox}{FIRST TAKE (5 min)}
Whose report would you send, Kai's or Rosa's? Cite one number or dotplot feature.\par
\answer{Not graded. Preserve the split between computation and choice.}
\end{firsttakebox}

\textit{Rules callout ``CENTER, SPREAD, AND OUTLIERS (keep on the page)'' is printed on the student sheet.}\par\medskip

\dokbadge{1}~\textbf{(a)} State the median and maximum weekly paid-work hours.\par
\answer{The median is 8 hours and the maximum is 30 hours.}

\dokbadge{2}~\textbf{(b)} Using the AP quartile convention, calculate $Q_1$, $Q_3$, the IQR, and both 1.5-IQR fences; identify any outlier. Then verify the mean and sample standard deviation. Show enough work to distinguish the methods.\par
\answer{Exclude the overall median. The seven values below it give $Q_1=4$; the seven above give $Q_3=12$. Thus IQR $=8$, the fences are $-8$ and $24$, and 30 is a high outlier. Also, $\sum x=126$, so $\bar{x}=8.4$; with squared-deviation sum $795.6$, $s_x=\sqrt{795.6/14}=7.538\ldots\approx7.5$ hours.}

\dokbadge{3}~\textbf{(c)} Adjudicate both reports. Whose IQR follows the AP convention, and what did the other student do differently? Choose mean and SD or median and IQR for the principal, cite the specific dotplot feature controlling your choice, and explain what the alternative pair alone could lead the principal to believe.\par
\answer{Kai's IQR of 8 follows the AP convention. Rosa included the overall median in both halves, giving $Q_1=4.5$, $Q_3=11$, and IQR $=6.5$. Report the median of 8 and corrected IQR of 8: the isolated 30 follows a 15-hour gap and exceeds the upper fence 24. Mean 8.4 and SD 7.5 alone could make outlier-driven spread seem routine, hiding that 14 of 15 students worked at most 15 hours.}

\begin{summaryexitbox}{EXIT}
One thing the video changed about my first take: \hrulefill
\end{summaryexitbox}

\end{document}
